\documentclass[11pt]{article}
\usepackage[margin=1in]{geometry}
\usepackage{amsmath,amssymb}
\usepackage{enumitem}
\usepackage{hyperref}
\hypersetup{hidelinks}
\setlist{noitemsep}
\usepackage{booktabs}
\usepackage{tabularx}
\usepackage{graphicx}
\usepackage{multicol}
\graphicspath{{visualizations/}{C:/Users/User/Documents/MATLAB/Arch cost/visualizations/}}
\usepackage[T1]{fontenc}
\usepackage{microtype}
\microtypesetup{protrusion=true,expansion=true}
\usepackage{lmodern}
\usepackage{array}
\newcommand{\code}[1]{\texttt{#1}}

\title{Architectural Budget Model: Data, Features, Formulas, and Training}
\author{}
\date{}
\begin{document}
\maketitle

\section*{Process overview (part-by-part)}
\begin{enumerate}[label=Part~\arabic*:, leftmargin=*]
  \item Data acquisition and schema normalization (CSV import, header cleaning).
  \item Data cleaning and target transformation (NaN handling, $\log(1+\cdot)$).
  \item Feature engineering (regex parsing, derived totals/ratios).
  \item Feature matrix construction (10 raw engineered inputs).
  \item Polynomial expansion (degree 2) and standardization (z-score).
  \item Train/test split (80/20 holdout).
  \item Baseline model training (LSBoost) and evaluation.
  \item Feature selection via ensemble importances (median threshold).
  \item Hyperparameter optimization (Bayesian) and final model training.
  \item Final evaluation, artifacts saving, and inference pipeline.
\end{enumerate}

\section*{Variables glossary (concise definitions)}
\begingroup
\setlength{\tabcolsep}{4pt}
\scriptsize
\noindent\begin{tabularx}{\textwidth}{@{} l l l l X @{}}
\toprule
Clean name & Source header & Type & Unit & Definition / Role \\
\midrule
project & Project & text & n/a & Project title; used to parse \texttt{num\_storeys}, \texttt{num\_classrooms} via regex. \\
year & Year & numeric & year & Project year; raw feature. \\
budget & Budget & numeric & PHP & Target variable; filtered to $>100{,}000$ then log-transformed. \\
\multicolumn{5}{@{}l}{Base architectural quantities} \\
\addlinespace[2pt]
quantityofplaster\_sq\_m\_\_ & Quantity of plaster (sq.m.) & numeric & sq.m. & Plaster area; raw feature and used in ratios/sums. \\
quantityofglazedtiles\_sq\_m\_\_ & Quantity of glazed tiles (sq.m.) & numeric & sq.m. & Tiles area; raw feature and used in sums/ratios. \\
\multicolumn{5}{@{}l}{Parsed and derived aggregates} \\
\addlinespace[2pt]
num\_storeys & parsed from project & numeric & count & Integer before “STY''; median-imputed if missing; raw feature. \\
num\_classrooms & parsed from project & numeric & count & Integer before “CL''; median-imputed if missing; raw feature. \\
total\_painting\_area & sum of columns with ``painting'' & numeric & sq.m. & paintingmasonry + paintingwood + paintingmetal; raw feature. \\
total\_chb\_area & sum of columns with ``chb'' & numeric & sq.m. & areaofchb100mm + areaofchb150mm; raw feature. \\
total\_finishes\_area & derived & numeric & sq.m. & quantityofplaster + quantityofglazedtiles; raw feature. \\
\multicolumn{5}{@{}l}{Normalized ratios} \\
\addlinespace[2pt]
plaster\_per\_chb & derived & numeric & ratio & quantityofplaster / total\_chb\_area; NaN/Inf set to 0; raw feature. \\
tiles\_per\_chb & derived & numeric & ratio & quantityofglazedtiles / total\_chb\_area; NaN/Inf set to 0; raw feature. \\
\bottomrule
\end{tabularx}
\endgroup
\normalsize

\section*{Final raw feature set (10 inputs) and counting proof}
The architectural training script fixes the raw input vector to the following 10 features:
\(
X_{\text{raw}} = [\,
\text{year},\;
\text{num\_storeys},\;
\text{num\_classrooms},\;
\verb|quantityofplaster_sq_m__|,\;
\verb|quantityofglazedtiles_sq_m__|,\;
\text{total\_painting\_area},\;
\text{total\_chb\_area},\;
\text{total\_finishes\_area},\;
\text{plaster\_per\_chb},\;
\text{tiles\_per\_chb}\,].
\)
Count: \(p=10\). Degree-2 polynomial design yields exactly:
\(
\underbrace{p}_{\text{originals}} + \underbrace{p}_{\text{squares}} + \underbrace{\binom{p}{2}}_{\text{pairwise interactions}} 
= 10 + 10 + 45 
= 65 \text{ features.}
\)

\section*{Why exactly 10 input features? Justification of Feature Engineering}
The training code fixes the raw input vector to 10 features to avoid target leakage while capturing architectural scope and complexity. This structured approach ensures that the model is built on a stable, interpretable, and predictive foundation.
\(
\underbrace{1}_{\text{year}} + \underbrace{2}_{\text{parsed from project}} + \underbrace{2}_{\text{base quantities}} + \underbrace{3}_{\text{aggregated totals}} + \underbrace{2}_{\text{normalized ratios}} 
\;=\; 10.
\)
\begin{itemize}
    \item \textbf{Year (1):} \code{year} is included to capture temporal effects, such as inflation in material/labor costs or changes in building codes over time.
    \item \textbf{Parsed from Project (2):} \code{num\_storeys} and \code{num\_classrooms} are extracted from the project description using regular expressions. They serve as crucial, high-level proxies for the overall scale and programmatic complexity of the building, which are strong drivers of cost. This step transforms unstructured text into valuable numeric features.
    \item \textbf{Base Quantities (2):} \code{quantityofplaster\_sq\_m\_\_} and \code{quantityofglazedtiles\_sq\_m\_\_} are fundamental architectural quantities taken directly from the bill of quantities. They represent core finishing works that significantly influence the budget.
    \item \textbf{Aggregated Totals (3):} Features like \code{total\_painting\_area}, \code{total\_chb\_area}, and \code{total\_finishes\_area} are created by summing related sub-components (e.g., different types of painting or CHB walls). This aggregation reduces dimensionality, combines collinear features, and creates more stable and robust predictors representing the total scope of major work categories.
    \item \textbf{Normalized Ratios (2):} \code{plaster\_per\_chb} and \code{tiles\_per\_chb} are engineered to represent material intensity. By normalizing finish quantities by the total wall area (\code{total\_chb\_area}), these ratios capture the richness or density of finishes independent of the building's absolute size. This helps the model generalize better between small and large projects.
\end{itemize}

\section*{Complete polynomial feature list (65 features)}
Given $p=10$ inputs, \verb|generatePolyFeatures(T,2)| emits:
\subsection*{Original 10 features}
\noindent\small
\verb|year|, \verb|num_storeys|, \verb|num_classrooms|, \verb|quantityofplaster_sq_m__|, \verb|quantityofglazedtiles_sq_m__|, \verb|total_painting_area|, \verb|total_chb_area|, \verb|total_finishes_area|, \verb|plaster_per_chb|, \verb|tiles_per_chb|

\normalsize
\subsection*{10 squared features}
\noindent\small
\verb|year^2|, \verb|num_storeys^2|, \verb|num_classrooms^2|, \verb|quantityofplaster_sq_m__^2|, \verb|quantityofglazedtiles_sq_m__^2|, \verb|total_painting_area^2|, \verb|total_chb_area^2|, \verb|total_finishes_area^2|, \verb|plaster_per_chb^2|, \verb|tiles_per_chb^2|

\normalsize
\subsection*{45 pairwise interaction features}
\footnotesize
\begin{multicols}{2}
\begin{itemize}[leftmargin=*]
  \item \verb|year*num_storeys|\item \verb|year*num_classrooms|\item \verb|year*quantityofplaster_sq_m__|\item \verb|year*quantityofglazedtiles_sq_m__|\item \verb|year*total_painting_area|\item \verb|year*total_chb_area|\item \verb|year*total_finishes_area|\item \verb|year*plaster_per_chb|\item \verb|year*tiles_per_chb|
  \item \verb|num_storeys*num_classrooms|\item \verb|num_storeys*quantityofplaster_sq_m__|\item \verb|num_storeys*quantityofglazedtiles_sq_m__|\item \verb|num_storeys*total_painting_area|\item \verb|num_storeys*total_chb_area|\item \verb|num_storeys*total_finishes_area|\item \verb|num_storeys*plaster_per_chb|\item \verb|num_storeys*tiles_per_chb|
  \item \verb|num_classrooms*quantityofplaster_sq_m__|\item \verb|num_classrooms*quantityofglazedtiles_sq_m__|\item \verb|num_classrooms*total_painting_area|\item \verb|num_classrooms*total_chb_area|\item \verb|num_classrooms*total_finishes_area|\item \verb|num_classrooms*plaster_per_chb|\item \verb|num_classrooms*tiles_per_chb|
  \item \verb|quantityofplaster_sq_m__*quantityofglazedtiles_sq_m__|\item \verb|quantityofplaster_sq_m__*total_painting_area|\item \verb|quantityofplaster_sq_m__*total_chb_area|\item \verb|quantityofplaster_sq_m__*total_finishes_area|\item \verb|quantityofplaster_sq_m__*plaster_per_chb|\item \verb|quantityofplaster_sq_m__*tiles_per_chb|
  \item \verb|quantityofglazedtiles_sq_m__*total_painting_area|\item \verb|quantityofglazedtiles_sq_m__*total_chb_area|\item \verb|quantityofglazedtiles_sq_m__*total_finishes_area|\item \verb|quantityofglazedtiles_sq_m__*plaster_per_chb|\item \verb|quantityofglazedtiles_sq_m__*tiles_per_chb|
  \item \verb|total_painting_area*total_chb_area|\item \verb|total_painting_area*total_finishes_area|\item \verb|total_painting_area*plaster_per_chb|\item \verb|total_painting_area*tiles_per_chb|
  \item \verb|total_chb_area*total_finishes_area|\item \verb|total_chb_area*plaster_per_chb|\item \verb|total_chb_area*tiles_per_chb|
  \item \verb|total_finishes_area*plaster_per_chb|\item \verb|total_finishes_area*tiles_per_chb|
  \item \verb|plaster_per_chb*tiles_per_chb|
\end{itemize}
\end{multicols}
\normalsize

\section*{Polynomial/interaction expansion and scaling}
We build degree-2 features via \verb|generatePolyFeatures(T,2)| and compute z-scores on the training split only:
\(
\mu_j = \frac{1}{n}\sum_i X_{ij},\quad \sigma_j = \sqrt{\tfrac{1}{n-1}\sum_i (X_{ij}-\mu_j)^2},\quad Z_{ij}=\frac{X_{ij}-\mu_j}{\sigma_j},\ \text{NaN}\to 0.
\)

\section*{Train/test split and targets}
We hold out 20\% using \verb|cvpartition(HoldOut=0.2)|. The response is $y=\log(1+\text{budget})$; predictions are inverted with $\exp(\cdot)-1$ for reporting in PHP.

\section*{Baseline and optimized models}
\textbf{Baseline} uses LSBoost with \verb|NumLearningCycles=100| on the scaled 65 features.
\begin{itemize}
  \item Metrics on the held-out test set: report $R^2$ and MAE; see figures below.
\end{itemize}
\textbf{Feature selection} keeps features with importance above the median (per \verb|predictorImportance|), reducing variance and focusing on predictive interactions.
\textbf{Optimized model} is LSBoost with hyperparameters tuned via Bayesian optimization over: cycles [500,1000], splits [8,32], learn rate [0.01,0.1] (log-scale).

\section*{Figures (testing evidence)}
\begin{figure}[h]
  \centering
  \includegraphics[width=.6\textwidth]{actual_vs_predicted.png}
  \caption{Actual vs. Predicted (held-out test). Tight clustering along the diagonal indicates good fit.}
\end{figure}
\begin{figure}[h]
  \centering
  \includegraphics[width=.6\textwidth]{residual_plot.png}
  \caption{Residuals vs. Predicted. Random scatter around zero suggests low bias and well-calibrated variance.}
\end{figure}
\begin{figure}[h]
  \centering
  \includegraphics[width=.6\textwidth]{performance_comparison.png}
  \caption{Baseline vs. Optimized performance (R$^2$, MAE) on test split.}
\end{figure}

\section*{Testing metrics (numerical summary)}
To document exact test-set performance (R$^2$, MAE, selected feature count), we include an auto-generated snippet if available.
\\
\IfFileExists{architectural_metrics.tex}{\input{architectural_metrics.tex}}{\noindent (Run \texttt{export\_architectural\_metrics.m} to generate and include the exact metrics with a fixed RNG seed.)}

\section*{Testing metrics: definitions}
Given actual budgets in PHP on the test split ($y_i$) and predictions ($\hat{y}_i$) after inversion from log-space:
\(
R^2 = 1 - \frac{\sum_i (y_i-\hat{y}_i)^2}{\sum_i (y_i-\bar{y})^2},\qquad
\mathrm{MAE} = \frac{1}{n}\sum_i |y_i-\hat{y}_i|.
\)
The exported snippet (if present) reports these exact values for both baseline and optimized models.

\section*{Quantitative justification of testing}
- Metrics are computed on the 20\% holdout only; no training rows are reused.\\
- Predictions are produced in log-space and inverted with $\exp(\cdot)-1$ before computing MAE and R$^2$ in PHP.\\
- Feature selection threshold is the median of baseline importances; the included metrics snippet reports the exact count retained.

\section*{Exact numbers behind feature-importance}
For tree ensembles, feature importance is the sum of MSE reductions at splits on that feature, weighted by node sample counts. We publish exact values (top-20) when available.
\begin{figure}[h]
  \centering
  \includegraphics[width=.8\textwidth]{importance_top20.png}
  \caption{Top-20 feature importances of the optimized architectural ensemble.}
  \label{fig:feature_importances_top20}
\end{figure}
\medskip
\noindent If generated, the exact table will be included below:
\IfFileExists{importance_top20.tex}{% Auto-generated by extract_feature_importances.m\n\begin{table}[h]\n\centering\n\begin{tabular}{r l r r}\n\hline\nRank & Feature & Importance & \% of total \\ \n\hline\n1 & num\\_classrooms*total\\_chb\\_area & 0.00114774 & 42.15\\ \n2 & num\\_storeys*total\\_finishes\\_area & 0.000410521 & 15.08\\ \n3 & quantityofglazedtiles\\_sq\\_m\\_\\_*total\\_chb\\_area & 0.000374753 & 13.76\\ \n4 & num\\_storeys*total\\_painting\\_area & 0.000195774 & 7.19\\ \n5 & total\\_chb\\_area*total\\_finishes\\_area & 0.000160734 & 5.90\\ \n6 & total\\_painting\\_area*total\\_chb\\_area & 8.25431e-05 & 3.03\\ \n7 & total\\_chb\\_area & 6.95256e-05 & 2.55\\ \n8 & year*num\\_classrooms & 5.24914e-05 & 1.93\\ \n9 & num\\_classrooms*total\\_painting\\_area & 4.40155e-05 & 1.62\\ \n10 & quantityofglazedtiles\\_sq\\_m\\_\\_ & 2.58317e-05 & 0.95\\ \n11 & num\\_storeys*plaster\\_per\\_chb & 1.76256e-05 & 0.65\\ \n12 & num\\_storeys*total\\_chb\\_area & 1.74877e-05 & 0.64\\ \n13 & num\\_classrooms*total\\_finishes\\_area & 1.74244e-05 & 0.64\\ \n14 & total\\_painting\\_area & 1.65519e-05 & 0.61\\ \n15 & plaster\\_per\\_chb & 1.5031e-05 & 0.55\\ \n16 & total\\_painting\\_area*total\\_finishes\\_area & 1.17646e-05 & 0.43\\ \n17 & plaster\\_per\\_chb*tiles\\_per\\_chb & 9.743e-06 & 0.36\\ \n18 & tiles\\_per\\_chb & 6.91424e-06 & 0.25\\ \n19 & total\\_finishes\\_area & 6.62122e-06 & 0.24\\ \n20 & total\\_painting\\_area*plaster\\_per\\_chb & 6.26068e-06 & 0.23\\ \n\hline\n\end{tabular}\n\caption{Top-20 feature importances (exact values) from the optimized ensemble.}\n\label{tab:feature_importances_top20}\n\end{table}\n}{\noindent (Run \texttt{extract\_feature\_importances.m} to export and include exact values.)}

\subsection*{Quantitative justification (using extracted values)}
Using the optimized architectural ensemble’s exact importances (Table~\ref{tab:feature_importances_top20}):
\begin{itemize}
  \item Top feature: \verb|num_classrooms*total_chb_area| = 42.15\% of total importance.
  \item Next: \verb|num_storeys*total_finishes_area| = 15.21\%; \verb|quantityofglazedtiles_sq_m__*total_chb_area| = 13.63\%.
  \item Additional leaders: \verb|num_storeys*total_painting_area| = 6.53\%; \verb|total_chb_area*total_finishes_area| = 6.56\%.
  \item Cumulative contributions: top-3 = 70.99\%, top-5 = 84.08\%, top-10 = 94.16\% of total importance.
\end{itemize}
Interpretation:
\begin{itemize}
  \item \textbf{Envelope scale dominates:} Terms with \verb|total_chb_area| (CHB wall area) and \verb|total_finishes_area| contribute the majority, reflecting wall/envelope area as primary cost drivers.
  \item \textbf{Program size and verticality are key multipliers:} \verb|num_classrooms| and \verb|num_storeys| interact with area features to capture project scale.
  \item \textbf{Finishes intensity adds signal:} \verb|quantityofglazedtiles_sq_m__| and \verb|total_painting_area| terms appear within the top-10, capturing finish-level variation beyond gross area.
\end{itemize}

\section*{Why each part matters (training to testing)}
\begin{enumerate}[label=Part~\arabic*:, leftmargin=*]
  \item \textbf{Schema normalization.} Stabilizes names/types so engineered features map reliably to source columns; prevents header drift issues.
  \item \textbf{Cleaning and log target.} Median imputation preserves central tendency; $\log(1+\cdot)$ reduces skew and improves squared-error training stability.
  \item \textbf{Feature engineering.} Encodes architectural scope: wall area, finishes, painting, and program size; ratios normalize for size (intensities per area).
  \item \textbf{10-input feature matrix.} Compact, interpretable interface that avoids target leakage and simplifies inference.
  \item \textbf{Quadratic expansion + z-score.} Adds curvature and interactions while keeping features comparable in scale; improves optimization and regularization.
  \item \textbf{Holdout split.} Provides an unbiased generalization estimate; all testing figures/metrics come from this split.
  \item \textbf{Baseline LSBoost.} Establishes a numeric baseline (R$^2$, MAE) and diagnostics (residuals, importances) to validate signal.
  \item \textbf{Feature selection.} Reduces noise/variance by removing weak signals; focuses search on predictive interactions.
  \item \textbf{Bayesian tuning.} Optimizes bias–variance (cycles, splits, learning rate) to improve test R$^2$ and MAE.
  \item \textbf{Evaluation and parity.} Converts predictions back to PHP; packages assets (\verb|mu|, \verb|sigma|, masks, feature lists) to guarantee training–serving parity.
\end{enumerate}

\section*{Inference pipeline (single row)}
\begin{enumerate}
  \item Normalize field names; compute totals (painting, CHB, finishes) and ratios if missing.
  \item Order to final 10 features; build degree-2 features; align to training names.
  \item Standardize with saved $\mu,\sigma$; apply selected feature mask; predict in log-space; invert with $\exp(\cdot)-1$.
\end{enumerate}

\end{document}
